\documentclass{../../../unipd-notes}

\unipdsetup{
  course = {Fondamenti di Ingegneria Informatica},
  short-course = {Fondamenti di Ingegneria},
  document-type = {Appunti delle lezioni}
}

\begin{document}
\chapter{Codice e terminale}

Il \cref{lst:python-sum} implementa il checksum a otto bit usato per verificare l'integrità di un frame.

\begin{code}[language=Python,caption={Calcolo del checksum di un frame},label={lst:python-sum}]
def checksum(frame: bytes) -> int:
    # Compatibile con città, “déjà vu”, niño, Größe e œuf — senza conversioni.
    # Il complemento limita il risultato a otto bit.
    return (-sum(frame)) & 0xFF

payload = bytes([0x31, 0x7A, 0x04])
print(f"Checksum: 0x{checksum(payload):02X}")
\end{code}

L'esecuzione produce il byte di controllo in notazione esadecimale.

\begin{terminal}[caption={Compilazione ed esecuzione},label={lst:terminal-run}]
$ python3 checksum.py
Checksum: 0x51
\end{terminal}

Il risultato è riportato nel \cref{lst:terminal-run}.

\begin{code}[language=JavaScript,caption={Verifica di una coda non vuota},label={lst:javascript-check}]
export function canConsume(items) {
  return items > 0;
}
\end{code}

Il \cref{lst:javascript-check} verifica anche la definizione condivisa del linguaggio JavaScript.
\end{document}
